\uuid{9T4d}
\titre{ Polynôme minimal et réduction }

\niveau{L2}
\module{Algèbre}
\chapitre{Réduction d'endomorphisme}
\sousChapitre{Polynôme annulateur}
\cpge{mp-mpi.reduction-endomorphismes.b, pc.reduction.a, psi.reduction-endomorphismes.a, pt.reduction-endomorphismes.a, mp-mpi.reduction-endomorphismes.c, pc.reduction.b, psi.reduction-endomorphismes.b, pt.reduction-endomorphismes.b, mp-mpi.reduction-endomorphismes.g, bcpst2.valeurs-propres-vecteurs-propres.a, bcpst2.valeurs-propres-vecteurs-propres.b}
\theme{}
\auteur{Quentin Liard}
\datecreate{ 2026-03-05}
\organisation{AMSCC}
\difficulte{1}





\contenu{

\texte{
On considère la matrice

\[
A=
\begin{pmatrix}
4 & 0 & 0\\
0 & 4 & 1\\
0 & 0 & 4
\end{pmatrix}.
\]
}

\begin{enumerate}

\item \question{Calculer le polynôme caractéristique de $A$.}

\item \question{Déterminer les sous-espaces propres de $A$.}

\item \question{$A$ est-elle diagonalisable ?}

\item \question{Déterminer le polynôme minimal de $A$.}

\end{enumerate}

}